\documentclass[11pt]{article}
\usepackage[margin=1in]{geometry}
\usepackage{amsmath,amssymb,amsthm,mathtools}
\usepackage{microtype}
\usepackage[hidelinks]{hyperref}

\newcommand{\E}{\mathbb E}
\newtheorem{theorem}{Theorem}
\newtheorem{proposition}{Proposition}[section]
\newtheorem{lemma}[proposition]{Lemma}
\theoremstyle{definition}
\newtheorem{definition}[proposition]{Definition}

\title{The critical two-point function of lattice $\phi^4_2$}
\author{Jiwoon Park}
\date{\today}

\begin{document}
\maketitle

\begin{abstract}
In this article, we show that the critical two-point correlation function of the lattice $\phi^4_2$ model in the strong coupling regime
follows that of the Ising universality class.
\end{abstract}

\subsection*{AI statement}

This document is generated using ChatGPT 6.1 Sol.

\section{Main results}

We consider real spins \(\phi_v\) on the square lattice with single-site
probability measure \(\rho_g(\mathrm d\phi)\propto
e^{-g(\phi^2-1)^2}\mathrm d\phi\), \(g>0\). The interaction weight
in finite volume is \(\exp(\beta\sum_{\langle u,v\rangle}\phi_u\phi_v)\),
where the sum counts each unoriented nearest-neighbor edge once.
Write \(\langle\cdot\rangle_\beta^0\) and
\(\langle\cdot\rangle_\beta^+\) for the infinite-volume free and plus
states, respectively, and set
\[
 G_{\phi,\beta}(x)=\langle\phi_0\phi_x\rangle_\beta^0,\qquad
 m_+(\beta)=\langle\phi_0\rangle_\beta^+,\qquad
 \beta_c(g)=\inf\{\beta>0:m_+(\beta)>0\}.
\]
Let \(\mu_{\rm I}\) denote the unique infinite-volume nearest-neighbor Ising
state at \(\beta_{\rm I}=\tfrac12\log(1+\sqrt2)\), and put
\(G_{\rm I}(x)=\mu_{\rm I}(\sigma_0\sigma_x)\).
Use \(|x|\) for the Euclidean norm and \(|x|_\infty\) for the max norm.

\begin{theorem}\label{thm:main}
There exist \(g_0<\infty\) and \(\omega>0\) such that for every fixed
\(g\ge g_0\) there is \(Z_\phi(g)>0\) with \(Z_\phi(g)=1+O(g^{-1})\)
and, uniformly as \(|x|\to\infty\),
\begin{equation}
 G_{\phi,\beta_c(g)}(x)
   = Z_\phi(g)^2G_{\rm I}(x)+O_g(|x|^{-1/4-\omega}).
 \label{eq:direct-bulk-phi-two}
\end{equation}
In particular \(G_{\phi,\beta_c(g)}(x)\sim C_g |x|^{-1/4}\) for a
constant \(C_g>0\).
\end{theorem}

The pure Ising inputs are the conditional-averaging estimates and
calibration stated in Lemma~\ref{lem:pure-reference}, from
\cite[Proposition 4.1 and Lemma 5.3]{BulkIsingRadialRG}.
Proposition~\ref{prop:microscopic-input} gives the starting lists for
the background and the physical insertions. Propositions
\ref{prop:blocking} and \ref{prop:linear} give the analytic and linear
blocking estimates. The massless tuning in Proposition~\ref{prop:tuning}
and the periodic comparison in Proposition~\ref{prop:periodic} then
supply the backgrounds used in Proposition~\ref{prop:two-point-comparison}.
The proof of Theorem~\ref{thm:main} uses that comparison to identify the
tuned temperature with \(\beta_c(g)\).

For the finite-volume argument, let \(\mathbb T_N=(\mathbb Z/N\mathbb Z)^2\),
with its nearest-neighbor edges. Use \(\langle\cdot\rangle_{N,\beta}\)
and \(\mu_{{\rm I},N}\) for the scalar-field and critical Ising
probability measures on this torus, and define
\[
 G_{\phi,\beta;N}(x)=\langle\phi_0\phi_x\rangle_{N,\beta},
 \qquad G_{{\rm I};N}(x)=\mu_{{\rm I},N}(\sigma_0\sigma_x).
\]
A fixed \(x\in\mathbb Z^2\setminus\{0\}\) is identified with its torus
image for all sufficiently large \(N\). For two commuting labels
\(t=(t_0,t_x)\), define the partition integral
\[
 \mathcal Z_{N,\beta}(t)=
 \int e^{\beta\sum_{\langle u,v\rangle}\phi_u\phi_v}
 (1+t_0\phi_0)(1+t_x\phi_x)
 \prod_{v\in\mathbb T_N}\rho_g(\mathrm d\phi_v).
\]
Write \(M_x=\{0,x\}\) and \(t^S=\prod_{u\in S}t_u\) for \(S\subseteq M_x\).
Source identities are in the ring
\(\mathbb C[t_0,t_x]/(t_0^2,t_x^2)\).
We write \(\varepsilon=g^{-1}\).
Positive constants \(c,C\) may change from line to line.
Constants in the RG bounds are uniform in the volume, scale index,
separation and sufficiently large \(g\), unless dependence is indicated.
The blocking ratio and starting radius are fixed before choosing \(g\).

\section{Microscopic interaction and sources}
\label{sec:amplitudes}

Use \(\Lambda\) for \(\mathbb Z^2\) or \(\mathbb T_N\). In torus
statements with two marks assume \(N>4|x|\) and \(N\ge5\).
Write \(d_\Lambda\) for graph distance and \(E(\Lambda)\) for the
nearest-neighbor edges.

\begin{definition}[Lists and norms]
\label{def:list-norms}
Let \(\Lambda=\mathbb Z^2\) or \(\mathbb T_N\), with the lattice
max-norm distance \(d_\infty\) (the periodic distance on
\(\mathbb T_N\)). All carrier diameters use \(d_\infty\).
Let \(B_r=B_r(0)\) be the closed ball of integer radius \(r\ge1\),
write \(B_r(z)\) for the ball centered at \(z\) and \(A+B_r\) for
the \(r\)-neighborhood of \(A\), and set
\(n_r(A)=|A+B_r|/r^2\). A list \(h=(h_{A,z})\) is indexed by
finite nearest-neighbor connected carriers \(A\subset\Lambda\)
and anchors \(z\in A\). The function \(h_{A,z}\) depends on
spins in \(A\), possibly on a proper subset, and
\(\|h_{A,z}\|_\infty=\sup_{\sigma\in\{-1,1\}^A}|h_{A,z}(\sigma)|\).
Set
\[
 \|h\|_r=\sup_{y\in\Lambda}
    \sum_{\substack{A\cap B_r(y)\ne\varnothing\\z\in A}}
       e^{n_r(A)}\|h_{A,z}\|_\infty.
\]
For a marked list \(U\) with fixed nonempty label set \(S\), each
carrier contains the sites of all labels in \(S\), and its anchor
is the site of one prescribed label of \(S\). We use the norm
\[
 \|U\|_{r,\bullet}=\sum_{A,z}e^{n_r(A)}\|U_{A,z}\|_\infty .
\]
In each case the space of lists with finite norm is a Banach space.
For a list whose coefficient functions have absolutely summable
sup norms, write \(\operatorname{ev}(h)=\sum_{A,z}h_{A,z}\).
This applies to marked lists on either lattice and to all lists
of finite norm on a finite torus. The projection \(\pi_A\)
on a carrier \(A\), abbreviated \(\pi\), subtracts the mean
over independent uniform spins on \(A\). A list is canonical if
each coefficient has zero such mean. We store the subtracted
constants separately. Equivariance under a lattice symmetry means
that applying it to the indices and to the spin configuration
preserves each coefficient value. A homogeneous list is
translation-equivariant. Spin-flip parity is taken coefficientwise.
\end{definition}

For \(a>0\), let
\(\rho_g^+\) have density proportional to \(e^{-g(a^2-1)^2}\).
Under the product of the single-site laws \(\rho_g\),
the representation \(\phi_u=a_u\sigma_u\) gives independent
amplitudes of law \(\rho_g^+\) and independent uniform signs,
with the two families independent. Denote expectation over the
amplitudes by \(\E_{\rm iid}\) and put \(\delta_u=a_u-1\).
For \(e=\{u,v\}\) write \(\sigma_e=\sigma_u\sigma_v\).

Fix a bounded open complex disk \(\mathcal B\) about \(\beta_{\rm I}\).

\begin{lemma}[Amplitude and edge bounds]
\label{lem:amplitude-moments}
For fixed \(p,C\ge0\) and all sufficiently large \(g\), one has
for \(\delta=a-1\)
\begin{equation}
 \E_{\rm iid}(|\delta|^p e^{Ca^2})\le C_{p,C}g^{-p/2},
 \qquad |\E_{\rm iid}\delta|\le C_0 g^{-1}.
 \label{eq:bulk-direct-well-absolute}
\end{equation}
For \(X_e=a_ua_v-1\) and \(d_e=e^{\beta X_e\sigma_e}-1\),
uniformly over \(\beta\in\mathcal B\), edges and signs, one has for
all sufficiently large \(g\)
\begin{equation}
 \|d_e\|_{L^p(\E_{\rm iid})}\le C_{p,\mathcal B}g^{-1/2}\quad(p\ge1),
 \qquad |\E_{\rm iid} d_e|\le C_{\mathcal B} g^{-1}.
 \label{eq:mic-edge-bound}
\end{equation}
\end{lemma}

\begin{proof}
The normalizing integral for \(\rho_g^+\) is at least \(c g^{-1/2}\)
by restriction to \([1,1+g^{-1/2}]\). Since
\((a^2-1)^2=(a+1)^2\delta^2\ge\delta^2\) and
\(e^{Ca^2}\le e^{2C+2C\delta^2}\), a Gaussian integral bounds
the numerator of the first estimate. Integration by parts on
\((0,\infty)\), using \(a e^{-g(a^2-1)^2}=0\) at both endpoints,
gives \(4g\,\E_{\rm iid}(a^4-a^2)=1\).
The identity \(a^4-a^2=2\delta+\delta^2(a^2+2a+2)\) and the first
estimate imply the signed bound.

Now \(X_e=\delta_u+\delta_v+\delta_u\delta_v\) and
\(|X_e|\le1+(a_u^2+a_v^2)/2\).
By independence and \eqref{eq:bulk-direct-well-absolute},
\[
 \E_{\rm iid}(|X_e|^p e^{C|X_e|})\le C_{p,C}g^{-p/2},
 \qquad
 \E_{\rm iid}X_e=2\E_{\rm iid}\delta+(\E_{\rm iid}\delta)^2.
\]
The bounds
\[
 |d_e|\le C_{\mathcal B}|X_e|e^{C_{\mathcal B}|X_e|},\qquad
 |d_e-\beta\sigma_e X_e|\le C_{\mathcal B} X_e^2 e^{C_{\mathcal B}|X_e|}
\]
then give \eqref{eq:mic-edge-bound}.
\end{proof}

In the square-free source ring,
\(1+t_u a_u\sigma_u=(1+t_u\sigma_u)(1+t_u\delta_u\sigma_u)\).
The first factor is independent of the amplitudes and equals
\(e^{t_u\sigma_u}\) in this ring. Define the torus amplitude polynomial
\begin{equation}
 W_N(\sigma;t)=\E_{\rm iid}\!\left[
       \prod_{e\in E(\mathbb T_N)}(1+d_e)
       \prod_{u\in M_x}(1+t_u\delta_u\sigma_u)\right].
 \label{eq:bulk-direct-amplitude-generating}
\end{equation}
For estimates we view \(t\) as ordinary complex variables on
\(|t_u|\le2\) and extract the square-free coefficients from the
analytic logarithm. For \(S\subseteq M_x\), write \([t^S]\)
for extraction of the coefficient of \(\prod_{u\in S}t_u\)
with all other label exponents zero.

\begin{proposition}[Microscopic interaction and sources]
\label{prop:microscopic-input}
For each fixed \(b\ge0\), all sufficiently large \(g\), and each
\(S\subseteq M_x\), there are coefficient functions
\(f^S_{A,z}(\beta)\), analytic in \(\beta\in\mathcal B\), indexed by finite
nearest-neighbor connected carriers \(A\subseteq\Lambda\) and
anchors \(z\in A\). For \(S\ne\varnothing\), all contributing
carriers contain \(S\) and are anchored at its prescribed site
(\(u\) for \(S=\{u\}\) and \(0\) for \(S=M_x\)). On the torus
\([t^S]\log W_N=\sum_{A,z} f^S_{A,z}\), and the plane coefficients
are given by the same local cluster formula. The bounds
\begin{equation}
 \begin{aligned}
 \sup_{y\in\Lambda}\sum_{\substack{A\ni y\\z\in A}}
     e^{b|A|}\|f^\varnothing_{A,z}\|_\infty&\le C_b\varepsilon,\\
 \sum_{A,z} e^{b|A|}\|f^S_{A,z}\|_\infty
    &\le C_b\varepsilon e^{-d_\Lambda(0,x)\mathbf 1_{\{S=M_x\}}}
       \qquad(S\ne\varnothing)
 \end{aligned}
 \label{eq:bulk-direct-carrier}
\end{equation}
hold uniformly in \(\beta\), \(N\) and the separation. The background
list is homogeneous, equivariant under quarter-turns, and even under
spin flip.
The parity of \(f^S_{A,z}\) under spin flip is \((-1)^{|S|}\).

Use \(\pi_A\) to separate the constant of each coefficient function
from its canonical list term. Denote the canonical background by
\(h^{\rm mic}\), adding
the temperature interaction
\((\beta-\beta_{\rm I})\sum_e\sigma_e\), with half of each edge term
assigned to each endpoint. Its site-rooted norm with weight
\(e^{b|A|}\) is at most \(C_b(\varepsilon+|\beta-\beta_{\rm I}|)\).
Denote the canonical marked lists by
\(U_{\rm mic}^S\), including the bare spin term \(\sigma_u\) for
\(S=\{u\}\), and the extracted scalars by \(c_{\rm mic}^S\).
Then \(c_{\rm mic}^{\{u\}}=0\). In the marked norms with
weight \(e^{b|A|}\), the singleton correction has size
\(O(\varepsilon)\), and the mixed list and scalar together have
size \(O(\varepsilon e^{-d_\Lambda(0,x)})\).
For the torus let \(c_{\rm mic}^{\varnothing}\) be the scalar from
the unmarked list and \(\mathcal Z_{{\rm I},N}\) the pure Ising
partition function with product uniform signs as base measure.
The exact partition identity is
\begin{equation}
 \frac{\mathcal Z_{N,\beta}(t)}{\mathcal Z_{{\rm I},N}}
 =
 e^{c_{\rm mic}^{\varnothing}+t^{M_x}c_{\rm mic}^{M_x}}
 \,\mu_{{\rm I},N}\!\left(
   \exp\!\left[\operatorname{ev}(h^{\rm mic})
         +\sum_{\varnothing\ne S\subseteq M_x}
                     t^S\operatorname{ev}(U_{\rm mic}^S)\right]\right).
 \label{eq:mic-partition}
\end{equation}
\end{proposition}

\begin{proof}
Make a graph on the factors in \eqref{eq:bulk-direct-amplitude-generating},
with an edge between factors when their site supports intersect.
Use \(d_i=d_e\) for an edge factor and \(d_i=t_u\delta_u\sigma_u\)
for a marked factor. A polymer \(P\) is a nonempty connected set
of factor vertices, \(A(P)\) is the union of its site supports, and
\(\zeta(P)=\E_{\rm iid}\prod_{i\in P}d_i\).
The components of any subset of factor vertices use disjoint
families of amplitudes. Expansion of \eqref{eq:bulk-direct-amplitude-generating}
therefore gives the polymer partition function with activities
\(\zeta(P)\) and with incompatibility \(A(P)\cap A(Q)\ne\varnothing\).

The degree of the factor graph is uniformly bounded, so it admits
a proper coloring with a fixed number \(q\) of colors. Apply
H\"older with exponent \(q\) to the products on the color classes.
Independence of the amplitudes in each class gives
\[
 |\zeta(P)|\le \prod_{i\in P}\|d_i\|_{L^q}
             \le (C_q g^{-1/2})^{|P|}.
\]
For \(|P|=1\) the signed mean bounds improve this to \(C\varepsilon\),
including for a marked factor. The number of connected \(m\)-factor
sets with \(y\in A(P)\) is at most \(C_1^m\), by traversing a
spanning tree from a factor incident to \(y\). Also \(|A(P)|\le2m\).
Let \(a(P)\) be the supremum of \(|\zeta(P)|\) over the signs,
\(\beta\in\mathcal B\) and \(|t_u|\le2\). The preceding bounds give
\[
 \sup_y\sum_{P:\,y\in A(P)} a(P)e^{(b+1)|A(P)|}\le C_b\varepsilon.
\]
Thus for each polymer \(Q\) the same weighted sum over polymers
incompatible with \(Q\) is bounded by \(|A(Q)|C_b\varepsilon\).
For \(C_b\varepsilon\le1\), this verifies the cluster criterion
for activity majorants \(a(P)e^{b|A(P)|}\) and charge \(|A(P)|\).

By the pinned cluster estimate \cite[Theorem 1]{Ueltschi}, the
absolute sum of connected logarithm terms whose generating union
contains \(y\), weighted by \(e^{b|\bigcup_i A(P_i)|}\), is at most
\(\sum_{P:\,y\in A(P)}a(P)e^{(b+1)|A(P)|}\).
Indeed one pins a polymer containing \(y\) and uses
\(e^{b|\bigcup_i A(P_i)|}\le\prod_i e^{b|A(P_i)|}\).
Connectedness of the incompatibility graph makes the full
generating union \(A=\bigcup_i A(P_i)\) a connected site set.
We keep \(A\) as carrier even when some of its spins do not enter
the resulting coefficient. The pinned estimate on the plane
follows by exhaustion, since it is uniform for finite polymer
families. It defines each plane coefficient by an absolutely
convergent local sum. The corresponding torus series is the
analytic logarithm of \(W_N\) obtained from the cluster expansion.

Divide each unmarked cluster contribution equally among the sites
of \(A\). This preserves the bound and the lattice symmetries.
For \(S\ne\varnothing\), Cauchy extraction of \([t^S]\) costs at most
one on the fixed polydisk. Any cluster with a nonzero extracted
coefficient has a marked factor at each \(u\in S\), so it can be
pinned at the prescribed marked site. When \(S=M_x\) its generating
union satisfies \(|A|\ge d_\Lambda(0,x)+1\). Using the pinned
estimate with \(b+1\) in place of \(b\) gives the additional
exponential in \eqref{eq:bulk-direct-carrier}. Simultaneous reversal
of the signs and the source labels leaves \(W_N\) invariant and
gives the asserted parity, also termwise in the local formula.

Subtracting the independent uniform-spin mean on each carrier costs
at most a factor two in the list norm. The absolute sum of the
removed marked constants is bounded by the same generating-carrier
estimate. Odd singleton terms have zero uniform-spin mean. Finally,
each original edge weight factors into
\(e^{\beta_{\rm I}\sigma_e}e^{(\beta-\beta_{\rm I})\sigma_e}(1+d_e)\).
Combining these identities with \(1+t_u\sigma_u=e^{t_u\sigma_u}\)
and the amplitude expectation gives \eqref{eq:mic-partition}.
\end{proof}

Since \(n_1(A)\le9|A|\), a bound on the site-rooted sum with weight
\(e^{b|A|}\) as in \eqref{eq:bulk-direct-carrier}, with \(b>9\),
bounds the unmarked \(r=1\) norm and the \(r=1\) total norm of
each marked list. For the unmarked norm, a root in \(B_1(y)\)
costs at most nine site-rooted sums.

\section{The analytic blocking map}
\label{sec:blocking}

We use the following spaces and local transformation from
\cite[Sections 5.1--5.2]{BulkIsingRadialRG}. We give the
estimates with the full carrier weight because the microscopic
lists need not have finite range.

\begin{definition}[One blocking step]
\label{def:blocking-map}
Let \(R=Lr\) with \(L\) an integer, and let
\(M=\lceil(\log L)^2\rceil\). A carrier is short if \(n_r(A)\le M\).
For sufficiently large \(L\), a short carrier is contained in
the interior of the cutting square \(B_R(z)\) about any
\(z\in A\). Denote by \(P_{R,z}\) the pure critical Ising
conditional expectation given the spins at sites \(y\) with
\(d_\infty(y,z)\ge R\). For tests supported in \(B_R(z)\),
use the finite-volume kernel in \(\{y:d_\infty(y,z)<R\}\)
with the prescribed spins on \(\{y:d_\infty(y,z)=R\}\).
This specifies the kernel at every boundary configuration and
agrees in the plane and on a torus whenever the square lifts
isometrically.

For finite input lists, expand \(\exp(\operatorname{ev}(h))\)
in ordered occurrences of input list terms, with repetitions
allowed. For an occurrence of carrier \(A\), put \(V_R(A)=A+B_{2R}\).
Join two occurrences if their \(V_R\) neighborhoods intersect.
In each monomial apply \(P_{R,z}\) to every isolated short
occurrence, leaving all other factors as they are. Take the
logarithm of the transformed series and use \(\pi\) term by term.
A connected coefficient with at least two occurrences is
assigned to the generating carrier \(Q=\bigcup_i V_R(A_i)\).
Divide an unmarked coefficient equally among its input anchors,
counting multiplicity. For a single occurrence use carrier
\(B_R(z)\) if it was averaged, and \(A\) otherwise, with the
original anchor in either case. For its estimates retain the
auxiliary generating carrier \(Q=V_R(A)\). The resulting
unmarked list is \(\mathcal R_{r,R}(h)\); on the torus the sum
of removed constants is \(\mathcal C_{r,R}(h)\).

In the marked version use the square-free variables
\(t_\ell\), so \(t_\ell^2=0\) and \(t^S=\prod_{\ell\in S}t_\ell\).
Root each output of type \(S\) at its prescribed anchor.
For \(p=1,2\) and input lists \(U_i\) with nonempty pairwise
disjoint label sets \(S_i\), let
\(\mathcal D^{(p)}_{r,R}(h)[U_1,\dots,U_p]\) be the coefficient
of \(t^{\bigcup_i S_i}\) obtained with input
\(h+\sum_i t^{S_i}U_i\). This is the mixed \(p\)-input derivative
of the transformation and takes values in pairs
\((\text{output list},\text{new scalar})\), with projections
\(\operatorname{pr}_{\rm list}\) and \(\operatorname{pr}_{\rm sc}\).
On pairs use the norm
\(\|(U,c)\|_{R,\bullet}=\|U\|_{R,\bullet}+|c|\).
An input scalar is independent of spins and is carried outside
the map.
\end{definition}

\begin{lemma}[Geometry of the carrier weights]
\label{lem:carrier-geometry}
There are absolute \(c,C>0\) such that, for every connected
carrier \(A\),
\[
 n_r(A)\ge c(1+\operatorname{diam}(A)/r),\qquad
 n_R(V_R(A))\le C(1+n_r(A)/L).
\]
Moreover \(A+B_{4R}\) is covered by at most
\(CL^2n_r(A)\) balls of radius \(r\).
The same bounds hold with periodic distances for \(R\le\vartheta N\),
where \(\vartheta>0\) is an absolute constant.
\end{lemma}
\begin{proof}
Along a path joining two sites at distance \(\operatorname{diam}(A)\),
select successive sites at distances from the starting point
spaced by \(3r\). The disjoint \(r\)-balls centered there lie in
\(A+B_r\), which proves the first bound. In particular every
short carrier lies in \(B_{CMr}(z)\) for each \(z\in A\), and
\(CMr<R\) for \(L\) sufficiently large.

For the second bound, choose a maximal \(10R\)-separated subset
of \(A\). Its \(13R\)-balls cover \(A+B_{3R}\). If there is
more than one center, a path in \(A\) starting at each center
and heading toward another leaves the \(R\)-ball about the
first center. Its initial part up to that radius has an
\(r\)-neighborhood of size at least \(cRr\), by the same
disjoint-ball argument. These neighborhoods lie in disjoint
\(2R\)-balls and are contained in \(A+B_r\). The number of
centers is thus at most \(C(1+|A+B_r|/(Rr))\), proving the
bound. For the last assertion choose instead a maximal
\(r\)-separated subset of \(A\). The disjoint balls of radius
\(\lfloor r/2\rfloor\) bound the number of its centers by
\(Cn_r(A)\). The \(B_{4R+r}\) ball about each center costs
at most \(CL^2\) balls of radius \(r\).
All constructions use the local cardinality bounds for
balls, which hold on the stated periodic scales.
\end{proof}

\begin{proposition}[Blocking with a diameter reserve]
\label{prop:blocking}
Fix \(L\) sufficiently large. There are \(\eta_L,C_L,c>0\),
uniform in \(r\) and in tori with \(R\le\vartheta N\), such
that the blocking map is analytic for \(\|h\|_r<\eta_L\).
Writing \(T_{r,R}=\mathrm D\mathcal R_{r,R}(0)\), one has
\begin{equation}
 \|\mathcal R_{r,R}(h)-T_{r,R}h\|_R\le C_L\|h\|_r^2,
 \qquad
 \|\mathrm D\mathcal R_{r,R}(h)-T_{r,R}\|\le C_L\|h\|_r .
 \label{eq:blocking-quadratic}
\end{equation}
The linear map \(T_{r,R}\) averages short carriers, keeps long
carriers and projects by \(\pi\). The same linear rule applies
to \(T^\bullet_{r,R}=\operatorname{pr}_{\rm list}
\mathcal D^{(1)}_{r,R}(0)\), with marked anchors.

Superscript \(\ge D\) means the sum of the generating terms
with \(\operatorname{diam}(Q)\ge D\), retained before scalar
extraction. For \(D\ge0\), \(p\in\{1,2\}\) and
\(\mathbf U=(U_1,\dots,U_p)\), the bounds are
\begin{equation}
\begin{aligned}
 \|\mathcal R_{r,R}^{\ge D}(h)\|_R
    &\le C_L e^{-cD/R}\|h\|_r,\\
 \|\mathcal D^{(p),\ge D}_{r,R}(h)[\mathbf U]\|_{R,\bullet}
    &\le C_L e^{-cD/R}\prod_i\|U_i\|_{r,\bullet},\\
 \|(\mathcal D^{(p),\ge D}_{r,R}(h)-
       \mathcal D^{(p),\ge D}_{r,R}(0))[\mathbf U]\|_{R,\bullet}
    &\le C_L e^{-cD/R}\|h\|_r\prod_i\|U_i\|_{r,\bullet}.
\end{aligned}
\label{eq:blocking-diameter}
\end{equation}
The right side of the first bound also bounds the rooted sum
before projection.
The marked estimates include the new scalar.

The map preserves translation invariance, square symmetry and
spin parity in the corresponding sectors. On each torus its
exact identity is
\begin{equation}
 \mu_{{\rm I},N}\bigl(e^{\operatorname{ev}(h)}\bigr)
    =e^{\mathcal C_{r,R}(h)}
      \mu_{{\rm I},N}\bigl(e^{\operatorname{ev}(\mathcal R_{r,R}(h))}\bigr),
 \label{eq:blocking-partition}
\end{equation}
The identity also holds in the square-free algebra with the
corresponding marked coefficients and scalars. The proposition
applies at \(r=1,\ L=r_0\) for a fixed sufficiently large \(r_0\).
\end{proposition}
\begin{proof}
The interiors of the replaced cutting squares in a monomial
are disjoint, and the remaining factors are measurable outside
their union. The nearest-neighbor Markov property gives the
same integral before and after the replacements. The symmetry
statements follow from the local rule, the anchor averaging
and the uniform-spin projection. Terms whose generating carrier
lifts isometrically have the same local rule in plane and torus.

We estimate the logarithm before projection. For fixed labeled
occurrences \(1,\dots,m\), let \(\mathcal G\) be their overlap
graph. Use auxiliary variables \(s_i^2=0\). The transformed
coefficient for any subset of occurrences is a product over
the components of its induced graph. The transformed series in
the \(s_i\) is therefore the hard-core polynomial of connected
subsets of \(\mathcal G\). A singleton activity contains its
linear image and any larger connected activity is the original
product. The logarithm coefficient of \(\prod_i s_i\) is a
sum over partitions into connected blocks. The coefficient of
a partition is the alternating sum over connected edge sets of
its incompatibility graph. This follows by expanding the
hard-core constraint into edge factors and grouping a graph
by its connected components. The sum vanishes for a disconnected
graph and is bounded in absolute value by its number of
spanning trees. To see the bound, fix an edge ordering and
group the connected edge sets by their Kruskal spanning tree.
For each tree, an eligible extra edge is one larger than every
edge on its tree path. The edge sets in that class consist of
the tree and any subset of its eligible edges, so their
alternating sum has magnitude at most one.

For each block partition and quotient tree, choose one internal
tree in each block and one contact edge for each quotient edge,
using fixed deterministic choices. Their union is a spanning
tree of \(\mathcal G\), with edges designated internal or
external. These designations recover the partition and quotient
tree. Writing \(\tau(\mathcal G)\) for the number of spanning
trees, the absolute coefficient is consequently bounded by
\[
 2^{m-1}\tau(\mathcal G)\prod_{i=1}^m\|f_i\|_\infty ,
\]
where \(f_i\) is the function of occurrence \(i\).
In particular the logarithmic coefficient has connected
generating carrier \(Q\).

By Lemma~\ref{lem:carrier-geometry} and subadditivity of
neighborhood cardinality,
\[
 2n_R(Q)\le Cm+(C/L)\sum_i n_r(A_i).
\]
Take \(L\) large enough that \(C/L\le1/2\).
For an unmarked output rooted at a ball \(B_R(y)\), one
occurrence meets \(B_{3R}(y)\), which costs \(O(L^2)\) input
root balls. For a marked output, root the tree at an input
occurrence of one prescribed mark and sum it in its total
marked norm. A child of an occurrence of carrier \(A\) meets
\(A+B_{4R}\), so summing over the child costs at most
\(CL^2n_r(A)\) input root balls. Additional marked inputs
can be summed in their total marked norms. Inducting from
leaves toward the root gives \(k\) such covering factors at
a vertex with \(k\) children. The reserve
\(e^{n_r(A)/2}\) absorbs \(n_r(A)^k\) at cost \(2^k k!\).
Thus the tree bound with weight \(e^{2n_R(Q)}\) costs at most
\(C_L^m\prod_i k_i!\) times the product of the input norms.
The factorials count the orderings of children. The number of
ordered rooted shapes is at most \(4^{m-1}\) and there are at
most \(m!\) labelings, canceled by the exponential factorial.
This yields a bound \(C_L^m\) times the product of input norms
for the homogeneous degree-\(m\) contribution, including
the order-one rule.

Every actual output carrier is contained in \(Q\). Since \(Q\)
is connected, \(n_R(Q)\ge cD/R\) on the diameter restriction,
by Lemma~\ref{lem:carrier-geometry}. One factor \(e^{n_R(Q)}\)
therefore pays for the output weight and the other yields
\(e^{-cD/R}\). This estimate is made before \(\pi\), which costs
at most two in sup norm, and bounds extracted constants as well.
Distinguishing \(p\) inputs in degree \(m\) costs at most \(m^p\).
For a difference of derivatives at \(h\) and zero the degree
is at least \(p+1\), so there is at least one factor \(\|h\|_r\).
Summing the resulting series on a sufficiently small ball
proves \eqref{eq:blocking-quadratic} and \eqref{eq:blocking-diameter}.

For finite lists the monomial identity and the logarithmic
identity hold first as power-series identities. Absolute
convergence gives the analytic identities on the stated
ball. Rooted estimates define the plane coefficients locally
and permit removal of carrier truncations on each torus, where
the total sum is bounded by finitely many rooted sums. This
gives \eqref{eq:blocking-partition}, including its marked
coefficients. Only geometry and the reference Markov property
enter this proof, including for the preliminary step.
\end{proof}

\section{Linear contraction and thermal tuning}
\label{sec:tuning}

For integers \(1\le r<R\) and an anchor \(z\), the cutting square
and conditional kernel are those of Section~\ref{sec:blocking}.
If \(F\) is supported in the square about \(z\) of radius \(s<R\), define
\[
 d_{R,z}(F)=\mu^+_{4R,z}(F)-\mu_{\rm I}(F),
\]
where \(\mu^+_{4R,z}\) is the nearest-neighbor critical Ising law in
the lattice disk of radius \(4R\) about \(z\), with plus boundary
condition. Put \(e_z=\frac14\sum_{w:\,w\sim z}\sigma_z\sigma_w\).
Let \(v_r^{\rm th}\) denote the homogeneous list with one coefficient
\(r^{-1}\pi P_{r,z}e_z\) on \(B_r(z)\) at each \(z\). Let \(v_r^u\)
denote the singleton list anchored at \(u\) with coefficient
\(r^{1/8}P_{r,u}\sigma_u\) on \(B_r(u)\).

\begin{lemma}[Pure reference estimates]\label{lem:pure-reference}
For each sufficiently large fixed ratio \(R/s\) there is a lower
threshold on \(s\) such that the following estimates hold uniformly
in \(z\) and in the bounded function \(F\) on \(B_s(z)\):
\begin{equation}
 |d_{R,z}(F)|\le C(s/R)^d\|F\|_\infty,\qquad
 \|P_{R,z}F-\mu_{\rm I}(F)\|_\infty
   \le C|d_{R,z}(F)|+C(s/R)^p\|F\|_\infty .
 \label{eq:reference-averaging}
\end{equation}
The pairs \((d,p)\) are \((1,3)\) for spin-flip-even tests invariant
under right-angle rotations about \(z\), \((1/8,9/8)\) for
spin-flip-odd tests, and \((1,2)\) for arbitrary spin-flip-even tests.
The constants are uniform over sufficiently large fixed ratios.

For all sufficiently large \(r\), \(\|v_r^{\rm th}\|_r\) and
\(\|v_r^u\|_{r,\bullet}\) are bounded above and below by positive
constants. For sufficiently large \(R/r\),
\begin{align}
 |d_{R,z}(r^{-1}\pi P_{r,z}e_z)|
   &\asymp r^{-2}(r/R), \label{eq:thermal-calibration}\\
 |d_{R,u}(r^{1/8}P_{r,u}\sigma_u)|
   &\asymp (r/R)^{1/8}. \label{eq:spin-calibration}
\end{align}
\end{lemma}

\begin{proof}
Equation~\eqref{eq:reference-averaging} is the pure Ising estimate of
\cite[Proposition 4.1]{BulkIsingRadialRG}. It extends to complex \(F\)
by applying it to real and imaginary parts. The distinguished-field
bounds and calibrations are \cite[Lemma 5.3]{BulkIsingRadialRG}.
For the thermal norm, the single coefficient has size comparable
to \(r^{-2}\), with \(O(r^2)\) anchors in a rooted sum and weights
bounded above and below on \(B_r(z)\). The same weights in the
marked norm give the bound for \(v_r^u\).
\end{proof}

Use \(T_{r,R}\) and \(T^\bullet_{r,R}\) for the pure linear maps on
backgrounds and marked lists of Proposition~\ref{prop:blocking}.

\begin{proposition}[Linear estimates]\label{prop:linear}
Fix \(0<\rho<1\). There is a sufficiently large integer \(L\), and
then an integer \(r_0\), such that for \(r\ge r_0\), \(R=Lr\), the
following assertions hold.
\begin{enumerate}
\item In the homogeneous, right-angle-invariant, even background
space, \(T_{r,R}v_r^{\rm th}=Lv_R^{\rm th}\). There are real
bounded complements to the thermal lines such that the matrix
of \(T_{r,R}\) in the line and complement coordinates is
\(\left(\begin{smallmatrix}L&B^{\rm th}_{r,R}\\0&D^{\rm th}_{r,R}\end{smallmatrix}\right)\),
with \(\|B^{\rm th}_{r,R}\|\le C_L\) and
\(\|D^{\rm th}_{r,R}\|\le\rho\).
\item On singleton odd lists anchored at \(u\),
\(L^{1/8}T^\bullet_{r,R}v_r^u=v_R^u\). The corresponding matrix of
\(L^{1/8}T^\bullet_{r,R}\) is
\(\left(\begin{smallmatrix}1&B^{\rm sp}_{r,R}\\0&D^{\rm sp}_{r,R}\end{smallmatrix}\right)\)
with \(\|B^{\rm sp}_{r,R}\|\le C_L\) and \(\|D^{\rm sp}_{r,R}\|\le\rho\).
\item On singleton even lists, the unscaled \(T^\bullet_{r,R}\)
has norm at most \(\rho\). This last assertion also holds on
\(\mathbb T_N\) when \(R\le\vartheta N\) for a sufficiently small
fixed \(\vartheta>0\).
\end{enumerate}
The line and complement coordinates have norms uniformly equivalent
to the corresponding list norms.
\end{proposition}

\begin{proof}
Set \(M=\lceil(\log L)^2\rceil\). By the carrier geometry in
Lemma~\ref{lem:carrier-geometry}, all short terms at \(z\) are
supported on \(B_s(z)\) with \(s=CMr\) for an absolute integer \(C\).
For the unchanged long terms, the comparison
\(n_R(A)\le C+(C/L)n_r(A)\) gives an extra \(e^{-cM}\).
In the background norm a root square at scale \(R\) can be covered
by \(O(L^2)\) root squares at scale \(r\). The long operators
therefore cost \(CL^2e^{-cM}\) in the background norm and
\(Ce^{-cM}\) in the marked norm, before any normalization.

For a background list \(h\), sum the short coefficients anchored at
\(z\) into \(F_z\). The rooted norm counts all anchors in a radius-\(r\)
square. By translation invariance their total weighted masses are
equal, so \(\|F_z\|_\infty\le C r^{-2}\|h\|_r\).
The rotation and spin-flip symmetries pass to \(F_z\).
By \eqref{eq:reference-averaging} and \eqref{eq:thermal-calibration}
one may subtract \(\alpha r^{-1}\pi P_{r,z}e_z\) to cancel its
disk response, where \(\alpha\) is independent of \(z\) and
\(|\alpha|\le CM\|h\|_r\).
The adjusted function has norm at most \(C(1+M)r^{-2}\|h\|_r\).
The projection \(\pi\) kills constants. Using \(p=3\) in
\eqref{eq:reference-averaging}, its projected image has sup norm at
most \(C(M/L)^3(1+M)r^{-2}\|h\|_r\).
The \(O(R^2)\) output anchors have bounded weights on their squares,
giving a background list bound \(CM^4 L^{-1}\|h\|_r\).

For a singleton odd list \(U\), the short sum at \(u\) has norm at
most \(\|U\|_{r,\bullet}\). The cancellation now subtracts a multiple
of \(v_r^u\) of size at most \(CM^{1/8}\|U\|_{r,\bullet}\), by
\eqref{eq:spin-calibration}. With \(p=9/8\), multiplication by
\(L^{1/8}\) gives a marked bound
\[
 C L^{1/8}(M/L)^{9/8}(1+M^{1/8})\|U\|_{r,\bullet}
 \le C M^{5/4} L^{-1}\|U\|_{r,\bullet}.
\]
Each distinguished carrier is short. The tower property gives the
claimed multipliers on the two lines. The preceding cancellation
bounds hold for every input representative modulo its line, and
therefore bound the quotient operator norms after adding the long
terms. Choose a bounded real dual functional equal to one on each
distinguished vector and complexify it. The calibration guarantees
uniform bounds for the resulting projections. Taking the kernels
as complements gives the stated triangular matrices. The shear
bounds follow from boundedness of the full operators, and the
transverse bounds follow from the quotient estimates upon increasing \(L\).

For the unscaled even mark, \((d,p)=(1,2)\) in
\eqref{eq:reference-averaging} gives a short operator bound
\(CM/L\) after projecting off constants, without any response
subtraction. Adding the long bound proves contraction. A short
cutting square at \(R\le\vartheta N\) lifts isometrically, and its
conditional kernel is the same as in the plane by the nearest-neighbor
Markov property. Thus the same estimate holds on the torus. Choose
\(L\) first so the operator estimates are below \(\rho\), then \(r_0\)
to meet the thresholds in Lemma~\ref{lem:pure-reference}.
\end{proof}

Fix these integers and choose \(\rho<\gamma<1\). Let
\(r_j=r_0L^j\) for \(j\ge0\). The preliminary map is
\(\mathcal R_{1,r_0}\) and each subsequent step uses
\(\mathcal R_{r_j,r_{j+1}}\). The preliminary bounds follow from
Proposition~\ref{prop:blocking} with ratio \(r_0\).

\begin{proposition}[Massless tuning]\label{prop:tuning}
For all sufficiently large \(g\) there is
\(\beta_*(g)=\beta_{\rm I}+O(\varepsilon)\) such that the plane lists
obtained from the microscopic background at \(\beta=\beta_*(g)\) satisfy
\begin{equation}
 \|h_j\|_{r_j}\le C\varepsilon\gamma^j\qquad(j\ge0).
 \label{eq:decay-background}
\end{equation}
The construction selects a unique root of the decaying-graph
equation on an interval \(|\beta-\beta_{\rm I}|\le C_0\varepsilon\)
with \(C_0\) fixed.
\end{proposition}

\begin{proof}
Use the complements of Proposition~\ref{prop:linear} and write
\(h_j=\tau_jv_{r_j}^{\rm th}+w_j\). The recursion takes the form
\[
 \tau_{j+1}=L\tau_j+B_j^{\rm th}w_j+a_j(\tau_j,w_j),\qquad
 w_{j+1}=D_j^{\rm th}w_j+b_j(\tau_j,w_j).
\]
Here \(B_j^{\rm th}=B^{\rm th}_{r_j,r_{j+1}}\), and similarly for
\(D_j^{\rm th}\). The remainders vanish quadratically, and their
derivatives are bounded by \(C_L(|\tau_j|+\|w_j\|_{r_j})\), by
Proposition~\ref{prop:blocking}. Write
\(\Phi_{j,i}=D^{\rm th}_{j-1}\cdots D^{\rm th}_i\) for \(j>i\) and
\(\Phi_{i,i}=\mathrm{Id}\).
For given remainders \(a_i,b_i\) the decaying solution with prescribed
\(w_0\) is
\begin{align*}
 w_j&=\Phi_{j,0}w_0+\sum_{i<j}\Phi_{j,i+1}b_i,\\
 \tau_j&=-\sum_{i\ge j} L^{j-i-1}(B_i^{\rm th}w_i+a_i).
\end{align*}
In the norm \(\sup_j\gamma^{-j}(|\tau_j|+\|w_j\|_{r_j})\),
the solution operator is bounded by the geometric sums with ratios
\(\rho/\gamma\) and \(\gamma/L\). Substituting the nonlinear
remainders defines a map on the sequence space. On a ball of radius
\(A\|w_0\|_{r_0}\) it is a contraction into the ball for fixed
large \(A\) and sufficiently small \(\|w_0\|_{r_0}\), by the
quadratic and derivative bounds. It yields an analytic graph
\(\tau_0=\psi(w_0)\) and a trajectory of norm \(O(\|w_0\|_{r_0})\).
Analyticity follows from the locally uniform contraction iterates in
the complex sequence space. The graph and trajectory are real for
real \(w_0\).

Put \(\Delta\beta=\beta-\beta_{\rm I}\). The sum of the microscopic
temperature coefficients at anchor \(z\) is \(2\Delta\beta\,e_z\).
Each temperature carrier is short. They have the same output
carrier at \(z\), so their preliminary linear image there is
\(2\Delta\beta\,\pi P_{r_0,z}e_z\).
The amplitude contribution to the microscopic norm and to its
\(\beta\)-derivative is \(O(\varepsilon)\) on a smaller fixed disk,
by Proposition~\ref{prop:microscopic-input} and Cauchy's formula.
The preliminary operator estimate gives
\[
 h_0(\beta,g)=2r_0\Delta\beta\,v_{r_0}^{\rm th}
       +O_{\|\cdot\|_{r_0}}((\Delta\beta)^2+\varepsilon),\qquad
 \partial_\beta h_0=2r_0 v_{r_0}^{\rm th}
       +O_{\|\cdot\|_{r_0}}(|\Delta\beta|+\varepsilon).
\]
It follows that \(f(\beta)=\tau_0(\beta,g)-\psi(w_0(\beta,g))\)
satisfies \(f(\beta)=2r_0\Delta\beta+O((\Delta\beta)^2+\varepsilon)\) and
\(f'(\beta)=2r_0+O(|\Delta\beta|+\varepsilon)\). For a fixed sufficiently
large \(C_0\), \(f\) changes sign between \(\Delta\beta=\pm C_0\varepsilon\)
and has positive derivative on this interval. At its root
\(\|w_0\|_{r_0}=O(\varepsilon)\), so the decaying graph estimate
proves \eqref{eq:decay-background}.
\end{proof}

\section{Periodic backgrounds}
\label{sec:periodic}

Fix a small numerical \(\kappa>0\). The projection \(\Pi_N\)
of a plane list keeps carriers of max-norm diameter less than
\(\kappa N\) and projects their sites to \(\mathbb T_N\).
For a homogeneous list use one representative of each anchor
modulo \(N\). For a rooted marked list, the chosen lift of its marked
anchor fixes the lifts of the retained carriers. In either case the
projection is bounded in the list norms for \(r\le\vartheta N\)
with \(\vartheta\) sufficiently small, by lifting all retained
carriers meeting a fixed norm root box to the chart around that box.

\begin{proposition}[Periodic comparison]\label{prop:periodic}
Let \(H_j^N\) denote the torus background at \(\beta_*(g)\) after
the preliminary step and \(j\) subsequent steps. For a sufficiently
small \(\vartheta>0\), for all \(j\) with \(r_j\le\vartheta N\),
\begin{equation}
 E_j:=\|H_j^N-\Pi_Nh_j\|_{r_j}
 \le C\varepsilon\gamma^j e^{-cN/r_j},\qquad
 \|H_j^N\|_{r_j}\le C\varepsilon\gamma^j.
 \label{eq:periodic-background}
\end{equation}
\end{proposition}

\begin{proof}
At the microscopic scale, clusters whose carriers fit the projected
chart lift uniquely after their anchors are fixed. Their activities
and combinatorial factors coincide. The other connected carriers
have at least \(cN\) sites, so the size-weight reserve in
Proposition~\ref{prop:microscopic-input} bounds the discrepancy by
\(C\varepsilon e^{-cN}\) in the scale-one norm.
The temperature term has identical local coefficients. At the
preliminary step, the Lipschitz estimate transports this discrepancy.
All new terms that fail chart comparison have generating diameter
at least \(cN\), so the tail estimate of Proposition~\ref{prop:blocking}
gives \(E_0\le C\varepsilon e^{-cN/r_0}\).

At a subsequent step compare the projected plane output with the
torus output from the projected plane input. On a generating union
of diameter less than \(cN\), for \(c\) sufficiently small relative
to \(\kappa\), all input and output carriers are retained, the cutting
kernels agree, and the same uniform-spin means are subtracted.
The two coefficient constructions therefore agree on these terms.
The remaining terms, including winding unions, have generating
diameter at least \(cN\). Their summed norm is at most
\(C\|h_j\|_{r_j}e^{-cN/r_{j+1}}\). Changing the torus input
from \(\Pi_Nh_j\) to \(H_j^N\) costs \(K E_j\) as long as both lie
in the analytic ball. Thus
\[
 E_{j+1}\le K E_j+C\varepsilon\gamma^j e^{-cN/r_{j+1}} .
\]
For \(X=N/r_j\) direct iteration gives
\[
 \frac{E_j}{\varepsilon\gamma^j}
 \le C(K/\gamma)^j e^{-cX L^j}
       +\frac C\gamma\sum_{m=0}^{j-1}(K/\gamma)^m e^{-cX L^m}
 \le C' e^{-c' X}.
\]
The last inequality holds uniformly for \(X\ge\vartheta^{-1}\) since
\(L^m\) grows geometrically. Combining this estimate with
\eqref{eq:decay-background} ensures that the torus input stays in
the analytic ball for sufficiently large \(g\), closing the induction.
\end{proof}

\section{The two marked sources}
\label{sec:two-sources}

Fix the tuned inverse temperature \(\beta_*(g)\) from
Proposition~\ref{prop:tuning} and two distinct sites \(0,x\).
For \(\varnothing\ne S\subseteq M_x\), let \(U_j^S\)
and \(c_j^S\) be the canonical list and the accumulated scalar
in the coefficient of \(t^S\) in the plane flow at scale \(r_j\).
Normalize both by \(r_j^{|S|/8}\).
Thus the preliminary step includes multiplication by
\(r_0^{|S|/8}\), and the step from \(j\) to \(j+1\) includes
multiplication by \(L^{|S|/8}\).
Write \(U_{N,j}^S,c_{N,j}^S\) for the periodic counterparts.
For superscripts we abbreviate \(\{u\}\) by \(u\) and
\(\{0,x\}\) by \(0x\). For marked lists set
\(\mathcal D_j^{(p)}=\mathcal D_{r_j,r_{j+1}}^{(p)}\).
In the following estimates, the constants may depend on the
fixed blocking parameters.

\begin{lemma}[Periodic comparison of marked coefficients]
\label{lem:periodic-marks}
For each fixed \(x\) and \(j\), as \(N\to\infty\),
\[
 \sum_{\varnothing\ne S\subseteq M_x}
 \left(\|U_{N,j}^S-\Pi_N U_j^S\|_{r_j,\bullet}
              +|c_{N,j}^S-c_j^S|\right)\longrightarrow0 .
\]
Here \(\Pi_N\) uses the carrier cutoff of
Section~\ref{sec:periodic} and the fixed marked anchor.
The plane lists omitted by \(\Pi_N\) have marked norm tending
to zero for each such \(x,j\).
\end{lemma}

\begin{proof}
For fixed \(x\) and large \(N\), the marked sites have unique lifts
in any retained chart containing them and based at a prescribed
marked anchor. Microscopic coefficients agree on these charts by
lifting each amplitude cluster. The errors in the microscopic
marked list and scalar have norm tending to zero by the weighted
carrier bound of Proposition~\ref{prop:microscopic-input}.
Indeed the generating support of an exceptional cluster has
diameter at least \(cN\), and the weight is applied before
extracting its scalar.

Write \(\Delta_{j,N}\) for the sum in the lemma. The preliminary
map at \(r=1\), \(R=r_0\) has bounded operator norms by
Proposition~\ref{prop:blocking}. The microscopic background
discrepancy tends to zero, and exceptional preliminary test unions
have diameter at least \(cN\). The diameter estimate and the
microscopic marked convergence give \(\Delta_{0,N}\to0\).
For subsequent steps, first compare the maps on projected lists.
The kernels, the projections off constants and the coefficient rules agree
on test unions that lift into the chart. Every exceptional
connected test union has diameter at least \(cN\), provided
\(r_{j+1}\le\vartheta N\). The diameter estimate in
Proposition~\ref{prop:blocking} therefore bounds the corresponding
list and scalar errors. On bounded sets of marked inputs,
the one- and two-input operators and their background differences
give
\[
 \Delta_{j+1,N}\le
 C_j\bigl(\Delta_{j,N}+E_j+e^{-cN/r_{j+1}}\bigr),
 \qquad E_j=\|H_j^N-\Pi_N h_j\|_{r_j}.
\]
The marked norms are bounded for each fixed number of steps
by the same operator estimates, including the fixed normalization
factors. Proposition~\ref{prop:periodic} gives \(E_j\to0\) at
fixed \(j\), so the recurrence proves the comparison. The
omitted plane tails tend to zero by absolute summability
in the marked norm.
\end{proof}

Write \(v_j^u=v_{r_j}^u\) for the pure spin list of
Section~\ref{sec:tuning}.
Use \(\rho<\gamma<1\) from Section~\ref{sec:tuning} and fix
\(\theta\in(\gamma,1)\).

\begin{proposition}[One-source flow]
\label{prop:one-source}
There is a real
\(Z_\phi(g)=1+O(\varepsilon)\) such that
\begin{equation}
 U_j^u=Z_\phi(g)v_j^u+\eta_j^u,\qquad
 \|\eta_j^u\|_{r_j,\bullet}\le C\varepsilon\theta^j .
 \label{eq:bulk-direct-one-mark-flow}
\end{equation}
The factor \(Z_\phi(g)\) is independent of \(u\).
The one-source scalars vanish at every scale, on the plane and
the torus.
\end{proposition}

\begin{proof}
At zero background the preliminary one-input operator sends the
bare spin list to \(v_0^u\) after normalization. At the tuned
temperature the microscopic background norm is \(O(\varepsilon)\).
The physical source differs from the bare spin by \(O(\varepsilon)\)
in the microscopic marked norm by
Proposition~\ref{prop:microscopic-input}. The bounded one-input
operator and its \(O(\varepsilon)\) difference from the operator
at zero background give
\(\|U_0^u-v_0^u\|_{r_0,\bullet}\le C\varepsilon\).

Let \(A_j=L^{1/8}\operatorname{pr}_{\mathrm{list}}\mathcal D_j^{(1)}(0)\).
Write \(U_j^u=\alpha_jv_j^u+w_j^u\) in the spin-line
splitting from Proposition~\ref{prop:linear}.
In these coordinates \(A_j\) has blocks
\[
 A_j=
 \begin{pmatrix}1&B_j^{\rm sp}\\0&D_j^{\rm sp}\end{pmatrix},
 \qquad \|B_j^{\rm sp}\|\le C,\quad
 \|D_j^{\rm sp}\|\le\rho .
\]
The longitudinal multiplier is one by nested conditional
expectation. The true update is \(U_{j+1}^u=(A_j+J_j)U_j^u\),
with
\[
 \|J_j\|\le C\|h_j\|_{r_j}\le C\varepsilon\gamma^j
\]
by Propositions~\ref{prop:blocking} and \ref{prop:tuning}.
The pure products have norm at most \(C_0\) for every
consecutive interval of scales. The transverse block is bounded
by the corresponding power of \(\rho\), and the off-diagonal
block by \(C\sum_{m\ge0}\rho^m\). Expanding each perturbed
product by the positions of the \(J_i\) gives the bound
\(C_0\prod_{i\ge0}(1+C_0\|J_i\|)\).
In particular \(\|U_j^u\|_{r_j,\bullet}\le C\).

The transverse recurrence and the preliminary estimate give
\[
 \|w_j^u\|_{r_j,\bullet}
 \le C\varepsilon\rho^j+
     C\varepsilon\sum_{i=0}^{j-1}\rho^{j-1-i}\gamma^i
 \le C_\theta\varepsilon\theta^j.
\]
The longitudinal increments satisfy
\(|\alpha_{j+1}-\alpha_j|\le
C(\|w_j^u\|_{r_j,\bullet}+\varepsilon\gamma^j)\).
Hence \(\alpha_j\) converges and its tail is
\(O(\varepsilon\theta^j)\). Its limit is \(1+O(\varepsilon)\)
because \(\alpha_0=1+O(\varepsilon)\).
Setting the other label to zero commutes with the microscopic
integral and each blocking map. Thus the one-source flow depends
only on \(u\). Choose the complements by translation from one
site. Translation invariance then gives the same limit for
all \(u\). The uniform norm bound on \(v_j^u\) gives
\eqref{eq:bulk-direct-one-mark-flow}. Finally the background is
even, the one-source input is odd, and all kernels and the
constant projection commute with spin flip. This makes each
one-source scalar increment, including the microscopic one, zero.
\end{proof}

\begin{proposition}[Mixed coefficient up to separation]
\label{prop:mixed-source}
Choose \(b_0>0\) sufficiently small and suppose
\(b_0|x|\ge r_0\). Let \(k\) be the last index with
\(r_k\le b_0|x|\), so that
\begin{equation}
 b_0|x|/L<r_k\le b_0|x|.
 \label{eq:separation-scale}
\end{equation}
Then
\begin{equation}
 \|U_k^{0x}\|_{r_k,\bullet}+|c_k^{0x}|
 \le C\varepsilon\theta^k
 \label{eq:mixed-bound}
\end{equation}
with the \(\theta\) of Proposition~\ref{prop:one-source}.
\end{proposition}

\begin{proof}
Set \(D=|x|\) and \(B_j=\|U_j^{0x}\|_{r_j,\bullet}+|c_j^{0x}|\).
Take \(b_0\) so small that the protective neighborhoods of the two
pure single-spin carriers are disjoint at each step up to \(k\),
including the preliminary step. In a step ending at \(r_{j+1}\)
those neighborhoods have radii at most \(r_j+2r_{j+1}\).
Their disjointness follows from
\(\|x\|_\infty\ge D/\sqrt2\) and \(r_{j+1}\le b_0D\).

At the preliminary step the microscopic mixed coefficient and
its scalar have norm \(O(\varepsilon e^{-cD})\).
The zero-background bilinear operator on the two bare spins
vanishes by the disjoint neighborhoods. Every remaining
connected term contains a background occurrence or a microscopic
one-source correction and has a generating test union containing
both sites. Proposition~\ref{prop:blocking}, used with
\(r=1\) and \(R=r_0\), bounds its joint list and scalar norm
by \(C\varepsilon e^{-cD/r_0}\), with normalization absorbed
into \(C\). Consequently \(B_0\le C\varepsilon e^{-cD/r_0}\).

For \(j<k\) the coefficient rule in the square-free ring gives
\begin{equation}
 (U_{j+1}^{0x},c_{j+1}^{0x})=
 L^{1/4}\Big((0,c_j^{0x})
      +\mathcal D_j^{(1)}(h_j)[U_j^{0x}]
      +\mathcal D_j^{(2)}(h_j)[U_j^0,U_j^x]\Big).
 \label{eq:mixed-update}
\end{equation}
The two label orders in the quadratic Taylor term cancel its
factor \(1/2\).
The accumulated scalar is spin independent and is carried along
separately. The terms linear in \((U_j^{0x},c_j^{0x})\) are
bounded in joint norm by \(K B_j\), for a scale-independent \(K\).

Write \(Z=Z_\phi(g)\). By disjointness,
\(\mathcal D_j^{(2)}(0)[Zv_j^0,Zv_j^x]=0\), including its
scalar part. Subtracting it from the bilinear term in
\eqref{eq:mixed-update} gives
\begin{align*}
 &\big(\mathcal D_j^{(2)}(h_j)-\mathcal D_j^{(2)}(0)\big)
       [U_j^0,U_j^x] \\
 &\quad + \mathcal D_j^{(2)}(0)[\eta_j^0,U_j^x]
         + \mathcal D_j^{(2)}(0)[Zv_j^0,\eta_j^x].
\end{align*}
Each contributing generating union has max-norm diameter at
least \(D/\sqrt2\) because it contains both marks. The joint diameter bound controls the displayed
terms by
\[
 C e^{-cD/r_{j+1}}\bigl(\|h_j\|_{r_j}
               +\|\eta_j^0\|_{r_j,\bullet}
               +\|\eta_j^x\|_{r_j,\bullet}\bigr).
\]
Here the remaining one-source factors have uniformly bounded
norm. It follows that
\[
 B_{j+1}\le K B_j+
 C\varepsilon\theta^j e^{-cD/r_{j+1}}.
\]
Iterating and using \(r_k/r_{k-m}=L^m\) bounds the
inhomogeneous sum by
\[
 C\varepsilon\theta^k\sum_{m\ge0}
    (K/\theta)^m e^{-c(D/r_k)L^m}\le C'\varepsilon\theta^k.
\]
The initial term obeys the same bound since
\(K^k e^{-cD/r_0}/\theta^k=(K/\theta)^k
e^{-c(D/r_k)L^k}\) is uniformly bounded.
This proves \eqref{eq:mixed-bound}.
\end{proof}

Define the residual periodic probability law by
\[
 \nu_{N,j}(d\sigma)=
 \frac{\exp\bigl(\operatorname{ev}(H_j^N)(\sigma)\bigr)}
      {\mu_{{\rm I},N}\bigl[\exp(\operatorname{ev}(H_j^N))\bigr]}
 \,\mu_{{\rm I},N}(d\sigma).
\]
The real coefficients of \(H_j^N\) and finite volume make this a
positive probability law.

\begin{lemma}[Comparison of an unscaled even test]
\label{lem:terminal-test}
Let \(F\) be an even test carried by a connected set \(A\)
with specified anchor \(u\in A\) on \(\mathbb T_N\).
Suppose \(k\le s\) and \(r_s\le\vartheta N\).
There is a \(q<1\) independent of \(F,k,s,N\) such that
\[
 |\nu_{N,k}(F)-\mu_{{\rm I},N}(F)|
 \le C e^{n_{r_k}(A)}\|F\|_\infty
        \bigl(\varepsilon\gamma^k+q^{s-k}\bigr).
\]
\end{lemma}

\begin{proof}
Let \(c\) be the uniform-spin mean of \(F\) on \(A\).
Insert \(F-c\) at scale \(k\) as one unscaled auxiliary
marked list \(f_k\), and carry \(c\) as its scalar.
The list norm is at most
\(C_F=2e^{n_{r_k}(A)}\|F\|_\infty\).
At scale \(j\) the linear transfer and scalar increment are
the two components of
\(\mathcal D_j^{(1)}(H_j^N)\), computed on the torus.
Use \(f_j^0\) for the analogous pure lists, with
\(f_k^0=f_k\). The unscaled even pure list map contracts
by Proposition~\ref{prop:linear}. Both the list map and
the scalar functional differ from their pure values in
operator norm by at most \(C\varepsilon\gamma^j\), by
Propositions~\ref{prop:blocking} and \ref{prop:periodic}.
For sufficiently large \(g\), both list maps have norm
at most some fixed \(q<1\) at all the chart scales.
It follows that
\[
 \|f_j\|_{r_j,\bullet}+\|f_j^0\|_{r_j,\bullet}
 \le 2C_F q^{j-k},\qquad
 \|f_j-f_j^0\|_{r_j,\bullet}
 \le C C_F\varepsilon\gamma^k(j-k)q^{j-k-1}\quad(j>k).
\]
The second bound is obtained by telescoping the products of
the two maps at each position \(i\in[k,j-1]\).

Let \(\ell_j(H_j^N)\) denote the scalar functional of the
one-input operator at scale \(j\).
The exact partition identity, with one square-free label,
gives
\[
 \nu_{N,k}(F)=c+
 \sum_{j=k}^{s-1}\ell_j(H_j^N)[f_j]
 +\nu_{N,s}(\operatorname{ev}(f_s)).
\]
There is the corresponding identity with \(H_j^N=0\),
\(f_j^0\) and \(\mu_{{\rm I},N}\).
The initial scalar cancels. Boundedness of the scalar
functional and the preceding estimates bound the difference
of the scalar sums by
\[
 C\sum_{j=k}^{s-1}
   \left(\|f_j-f_j^0\|_{r_j,\bullet}
          +\varepsilon\gamma^j\|f_j\|_{r_j,\bullet}\right)
 \le C C_F\varepsilon\gamma^k.
\]
The two terminal expectations are bounded in absolute value
by their list norms, since the measures are positive.
This proves the claimed estimate.
\end{proof}

\begin{proposition}[Two-point comparison at the tuned temperature]
\label{prop:two-point-comparison}
With \(k\) as in Proposition~\ref{prop:mixed-source},
\begin{equation}
 \limsup_{N\to\infty}
 \left|r_k^{1/4}G_{\phi,\beta_*;N}(x)-
 Z_\phi(g)^2 r_k^{1/4}G_{{\rm I};N}(x)\right|
 \le C\varepsilon\theta^k.
 \label{eq:two-point-periodic-comparison}
\end{equation}
Here \(x\) and \(k\) are held fixed in the periodic limit.
\end{proposition}

\begin{proof}
The partition identity at scale \(k\), after canceling the
background scalar and using the vanishing one-source scalars,
is
\[
 \frac{\mathcal Z_{N,\beta_*}(r_k^{1/8}t)}
      {\mathcal Z_{N,\beta_*}(0)}
 = e^{t_0t_xc_{N,k}^{0x}}\,
    \nu_{N,k}\!\left[
     \exp\!\left(t_0\operatorname{ev}(U_{N,k}^0)
             +t_x\operatorname{ev}(U_{N,k}^x)
             +t_0t_x\operatorname{ev}(U_{N,k}^{0x})\right)\right].
\]
The argument \(r_k^{1/8}t\) multiplies both labels.
Taking the coefficient of \(t_0t_x\) yields
\begin{equation}
 \begin{split}
 r_k^{1/4}G_{\phi,\beta_*;N}(x)
 &=c_{N,k}^{0x}
   +\nu_{N,k}(\operatorname{ev}(U_{N,k}^{0x}))\\
 &\quad+\nu_{N,k}\bigl(
    \operatorname{ev}(U_{N,k}^0)\operatorname{ev}(U_{N,k}^x)\bigr).
 \end{split}
 \label{eq:direct-two-source-exact}
\end{equation}

Set \(F_k=\operatorname{ev}(v_k^0)\operatorname{ev}(v_k^x)\)
on the two lifted cutting squares, viewed on the torus.
The marked norm bounds the evaluation norm, and chart projection
does not increase it. Lemma~\ref{lem:periodic-marks} and
Propositions~\ref{prop:one-source} and \ref{prop:mixed-source}
therefore give
\[
 r_k^{1/4}G_{\phi,\beta_*;N}(x)
   =Z_\phi(g)^2\nu_{N,k}(F_k)
      +O(\varepsilon\theta^k)+o_N(1).
\]
The test \(F_k\) is even with uniformly bounded sup norm.
By \eqref{eq:separation-scale}, a square of radius \(C_{b_0}r_k\)
contains both carriers and has \(n_{r_k}\) bounded uniformly
in \(x\). Lemma~\ref{lem:terminal-test}, with the largest \(s\)
such that \(r_s\le\vartheta N\), now gives
\[
 \nu_{N,k}(F_k)=\mu_{{\rm I},N}(F_k)+
 O(\varepsilon\gamma^k)+o_N(1).
\]
Finally, under \(\mu_{{\rm I},N}\), condition on the exterior
of the two cutting interiors. The disjoint interiors have
conditionally independent spins and their conditional means are
the defining kernels of \(v_k^0\) and \(v_k^x\).
Consequently \(\mu_{{\rm I},N}(F_k)=r_k^{1/4}G_{{\rm I};N}(x)\)
for all sufficiently large \(N\). Since \(\gamma<\theta\),
this proves \eqref{eq:two-point-periodic-comparison}.
\end{proof}

\section{Thermodynamic limit and criticality}
\label{sec:criticality}

\begin{lemma}[Fixed-separation torus limits]\label{lem:thermodynamic-limits}
For each fixed \(\beta\ge0\), \(g>0\) and \(x\in\mathbb Z^2\),
\[
 \lim_{N\to\infty}G_{\phi,\beta;N}(x)=G_{\phi,\beta}(x),\qquad
 \lim_{N\to\infty}G_{{\rm I};N}(x)=G_{\rm I}(x).
\]
\end{lemma}

\begin{proof}
Let \(m_{p,N}=\langle|\phi_0|^p\rangle_{N,\beta}\). Integration by
parts in \(\phi_0\), using the quartic tails, gives
\[
 4g(m_{4,N}-m_{2,N})=
 1+\beta\sum_{z:\,z\sim0}\langle\phi_0\phi_z\rangle_{N,\beta}
 \le 1+4\beta m_{2,N}.
\]
Since \(m_{2,N}\le m_{4,N}^{1/2}\), the fourth moments are uniformly
bounded for fixed \(g,\beta\). The periodic extensions to \(\mathbb Z^2\)
are therefore locally tight. Any limit is stationary and spin-flip
symmetric.

For each finite set, the finite-volume conditional kernels coincide
with the infinite-lattice kernels on fixed cylinder tests when \(N\)
is sufficiently large. For a bounded continuous test of the interior
variables, the kernel is continuous in its boundary values by quartic
domination on compact sets. Passing the conditional identity to the
limit and then using a monotone class gives the DLR equations.
They imply that finite marginals are absolutely continuous. Moreover,
the lattice pointwise ergodic theorem applied to \(\phi_0^2\) gives
almost surely bounded centered box averages of \(\phi_v^2\).
Thus the limits meet the temperedness requirement of
\cite[Appendix D]{GPPSCurrents}. By
\cite[Theorem 1.2]{GPPSCurrents}, each stationary symmetric limit is the
symmetric mixture of the plus and minus states, which is also the
free state. The uniform fourth moment bound gives uniform
integrability of \(\phi_0\phi_x\), so the correlations converge to
those of that state.

For the pure Ising model, local subsequential limits are critical
nearest-neighbor Gibbs measures. Uniqueness at \(\beta_{\rm I}\)
\cite{DSTIsingUnique} gives the second assertion.
\end{proof}

\begin{proof}[Proof of Theorem~\ref{thm:main}]
Proposition~\ref{prop:two-point-comparison} and
Lemma~\ref{lem:thermodynamic-limits} imply, at \(\beta=\beta_*(g)\),
\[
 G_{\phi,\beta_*}(x)
 =Z_\phi(g)^2G_{\rm I}(x)+O(\varepsilon r_k^{-1/4}\theta^k).
\]
Here \(r_k\asymp |x|\) by \eqref{eq:separation-scale}, so for
\(\omega=-\log\theta/\log L>0\) the error is
\(O(\varepsilon |x|^{-1/4-\omega})\).
The pure reference has \(G_{\rm I}(x)\sim C_{\rm I}|x|^{-1/4}\)
with \(C_{\rm I}>0\) \cite{CIMBulkSpin}. Thus the correlation at
\(\beta_*(g)\) tends to zero and its lattice sum diverges.

The three critical inverse temperatures for magnetization, finite
susceptibility and long-range order coincide by
\cite[Proposition C.1]{GPPSCurrents}. For \(\beta<\beta_c(g)\),
the free two-point function is summable. For
\(\beta>\beta_c(g)\), it has strictly positive infimum over \(x\),
using its monotonicity in \(\beta\) by the Griffiths inequality.
The two properties at \(\beta_*(g)\) force
\(\beta_*(g)=\beta_c(g)\). The preceding comparison proves
\eqref{eq:direct-bulk-phi-two} with
\(C_g=C_{\rm I}Z_\phi(g)^2>0\).
\end{proof}

\begingroup
\small
\begin{thebibliography}{9}
\setlength{\itemsep}{2pt}

\bibitem{BulkIsingRadialRG}
OpenAI.
\emph{Conformal universality of bulk Ising correlations under weak interactions}.
Preprint, 2026.
\href{https://github.com/openai/math/blob/main/preprints/Conformal-universality-of-bulk-Ising-correlations-under-weak-interactions-September-23-2026/paper.pdf}{Manuscript}.

\bibitem{Ueltschi}
D.~Ueltschi.
\emph{Cluster expansions and correlation functions}.
\href{https://arxiv.org/abs/math-ph/0304003}{arXiv:math-ph/0304003}.

\bibitem{GPPSCurrents}
T.~S.~Gunaratnam, C.~Panagiotis, R.~Panis and F.~Severo.
\emph{Random tangled currents for \(\varphi^4\): translation invariant Gibbs measures and continuity of the phase transition}.
\href{https://arxiv.org/abs/2211.00319}{arXiv:2211.00319}.

\bibitem{DSTIsingUnique}
H.~Duminil-Copin, V.~Sidoravicius and V.~Tassion.
\emph{Continuity of the phase transition for planar random-cluster and Potts models with \(1\le q\le4\)}.
\href{https://arxiv.org/abs/1505.04159}{arXiv:1505.04159}.

\bibitem{CIMBulkSpin}
D.~Chelkak, K.~Izyurov and R.~Mahfouf.
\emph{Universality of spin correlations in the Ising model on isoradial graphs}.
\href{https://arxiv.org/abs/2104.12858}{arXiv:2104.12858}.

\end{thebibliography}
\endgroup
\end{document}
